\subsection{Materia necesaria}

\begin{definicion}{Revision continua versus periodo fijo}
En revision continua cada producto puede tener su propio \(Q^*\) y ROP. En periodo fijo conjunto, se aprovecha un despacho comun; por eso se busca un ciclo \(T\) compartido.
\end{definicion}

\begin{formulas}{Revision continua}
\[
Q_i^*=\sqrt{\frac{2SD_i}{h_i}},
\qquad
ROP_i=d_iL+z\sigma_i\sqrt L,
\qquad
N_i=\frac{D_i}{Q_i}.
\]
\end{formulas}

\begin{formulas}{Periodo fijo conjunto}
Si \(S\) se paga una vez por despacho conjunto:
\[
T^*=\sqrt{\frac{2S(365)}{\sum_i h_i d_i}}.
\]
Luego:
\[
Q_i=d_iT^*,
\qquad
R_i=d_i(T^*+L)+z\sigma_i\sqrt{T^*+L}.
\]
\end{formulas}
